\scanpage{148}
s'interprètent comme les hom (17) satisfaisant
la \uline{première relation} (18). \deleted{\unsure{C'est lorsqu|C'est à lorsqu}} C'est la
deuxième relation qui précise que l'on se restreint
aux connexions intégrables sur $P$.

Les \deleted{faisceau des} automorphismes de $P$ \uline{compatibles avec} $d_P$
\deleted{2) \uline{Stratifications sur torseurs affines}}
est \deleted{\unsure{form}} sous-gpe du gpe $\Gamma(S,E)$ des automorphismes
de $P$ formé des \deleted{$x\in\Gamma(S,E$} sections \uline{horizontales} de $E$,
i.e. $H_{\mathrm{DR}}^0(S,E)$.

Les descriptions qui précèdent sont \inserted{fonctorielles en $E$, et} compatibles
aux images inverses envisagés (sic) plus haut.

2) \uline{Stratifications sur torseurs affines}

Supposons donnée sur $E$ une stratification rel. à $T$.
Alors $E$ définit un faisceau \inserted{de modules $E_{\mathrm{strat}}$} du site \unsure{stratifiant|stratifié}
$(S/T)_{\mathrm{strat}}$, donc on peut former
\[
H^*_{\mathrm{strat}}(S,E)\overset{\unsure{\mathrm{dfn}|\mathrm{déf}}}{=}H^*((S/T)_{\mathrm{strat}},E_{\mathrm{strat}}). \tag*{(19)}
\]
\deleted{La stratification de $E$}
\deleted{On a défini dans [\quad] un hom naturel}

Une donnée de recollement \unsure{d'infinitésimale|infinitésimale} d'ordre $\infty$ sur
$E$ revient à la donnée d'un homomorphismes
\margin{On dit parfois que $\delta_E^\infty$ est la connexion.}
\[
\delta_E^\infty:E\longrightarrow P^\infty_{S/T}(E) \tag*{(20)}
\]
satisfaisant la condition
\[
\varepsilon_E\delta_E^\infty:E\longrightarrow P^\infty_{S/T}(E)
\xrightarrow{\varepsilon_E} E\quad\text{est l'identité.} \tag*{(21)}
\]
\margin{$P^\infty_{S/T}(E)^+$, idéal d'augmentation}
\inserted{On préfèrera $\deleted{\unsure{d_E}}d_E^+:E\longrightarrow P^\infty_{S/T}(E)^+,\quad
d_E^+(x)=d^\infty_{\deleted{\unsure{E}}}(x)-\delta_E^\infty(x)$}
\deleted{(21) La condition pour la donnée de recollement inf. soit
une stratification est une condition sur $\delta_E$ qui}
% slanted
\margin{plus bas, ds n° 3.}
% archive sheet number: 147
